% appA —— 附录 A（原文 wcf-paper.tex 3242–3392 行，PDF 51–55 页）
\section{验证若干 $SU(2)$ 规范理论的 KS 恒等式}
\label{app:ks-product}

有一个富有启发性的办法，可以证明涉及 $\CK_{1,0}$ 与 $\CK_{0,1}$ 的那个简单公式。考虑一列数 $x_n$，它们满足递推关系
\begin{equation} \label{eq:appA-five-recursion}
x_{n+1} x_{n-1} = 1 - x_n.
\end{equation}
令人惊讶的是，这个递推关系是周期的，周期为五：
\begin{equation} \label{eq:appA-five-period}
x_2 = \frac{1-x_1}{x_0}, \quad x_3 = \frac{x_0+ x_1-1}{x_0 x_1}, \quad x_4 = \frac{1-x_0}{x_1}, \quad x_5 = x_0, \quad x_6=x_1.
\end{equation}
现在令 $X_{1,0}=X_{1,0}^{(1)}= x_1^{-1}$，$X_{0,1}=X_{0,1}^{(1)}= x_0$。我们的策略是：取一列适当的 KS 变换 $\CK_n$，依次定义变换 $(X_{1,0}^{(n+1)},X_{0,1}^{(n+1)}) = \CK_n (X_{1,0}^{(n)},X_{0,1}^{(n)})$，直到在 $(X_{1,0}^{(1)},X_{0,1}^{(1)})\to (X_{1,0}^{(N)},X_{0,1}^{(N)})$ 处得到恒等变换为止（在第一个例子中 $N=6$，但在其余例子中将会是无穷大）。为避免记号累赘，下文中我们不再标出上标 ${}^{(n)}$。

若施加 $\CK_{1,0}$，它不改变 $X_{1,0} = x^{-1}_1$ 的值，而是把 $X_{0,1}$ 变为 $X_{0,1}=x_0(1-x^{-1}_1)^{-1} =- x_1 x^{-1}_2$。注意 $X_{1,1} = X_{1,0} X_{0,1} = -x_2^{-1}$。

接着可以施加 $\CK_{1,1}$：它使 $X_{1,1} = -x^{-1}_2$ 保持不变，而改变 $X_{1,0}$。结果现在 $X_{0,1}=- x_1 x^{-1}_2 (1-x^{-1}_2)^{-1}= x_3^{-1}$。

若施加 $\CK_{0,1}$，则 $X_{0,1} = x_3^{-1}$，而 $X_{1,0} = - x_3 x^{-1}_2 (1-x^{-1}_3)=x_4$。若施加 $\CK_{1,0}^{-1}$，则 $X_{1,0} = x_4$，$X_{0,1} =x^{-1}_3 (1-x_4) = x_5$。最后，若施加 $\CK_{0,1}^{-1}$，我们得到 $X_{0,1}=x_5=x_0$，$X_{1,0} = x_4 (1-x_5)^{-1} = x^{-1}_6=x^{-1}_1$。于是便导出了所要的
\begin{equation} \label{eq:appA-pentagon}
\CK_{0,1}^{-1}\CK_{1,0}^{-1}\CK_{0,1}\CK_{1,1}\CK_{1,0} = 1.
\end{equation}

对更有意思的公式而言，这是一个有用的热身练习。现在考虑另一个递推关系：
\begin{equation} \label{eq:appA-sq-recursion}
x_{n+1} x_{n-1} = (1- x_n)^2.
\end{equation}
这个递推关系一般不是周期的：它的通解为
\begin{equation} \label{eq:appA-sq-solution}
x_n = -\frac{\cosh^2 (a n +b)}{\sinh^2 a}.
\end{equation}
我们仍可把这个递推关系与 $\CK$ 因子的一个乘积联系起来。

我们仍从 $X_{0,1}=x_0$ 与 $X_{1,0} = x^{-1}_1$ 出发。若施加 $\CK_{1,0}^2$，结果是 $X_{1,0}=x^{-1}_1$，$X_{2,1} = x_1^{-2} x_0 (1- x^{-1}_1)^{-2}=x^{-1}_2$。若施加 $\CK_{2,1}^2$，结果是 $X_{2,1}=x^{-1}_2$，$X_{3,2}=x_2^{-2} x_1 (1-x^{-1}_2)^{-2} = x^{-1}_3$。我们可以对一切 $n$ 继续施加 $\CK_{n+1,n}^2$，循着递推关系走到任意大的 $n$。借助关系 $X_{n+1,n} = x^{-1}_{n+1}$ 与 $X_{n,n-1} = x^{-1}_{n}$ 在 $n\to\infty$ 时的极限，便可算出这个无穷乘积。若取 $a,b$ 的实部为正，则 $X_{1,1} = e^{-2 a}$，且 $X_{1,0} =- e^{-2 b} (1-e^{-2 a})^2$。

另一方面，我们也可以沿相反方向追随这个递推关系：$X_{0,1}=x_0$ 与 $X_{1,0} =x^{-1}_1$ 在 $\CK_{0,1}^2$ 之下变为 $X_{0,1}=x_0$ 与 $X_{1,0} = x^{-1}_1 (1-x_0)^2 = x_{-1}$。而 $\CK_{1,0}^2$ 把它送到 $X_{1,0} = x_{-1}$ 与 $X_{0,1} = x_0 (1-x_{-1})^{-2} = x^{-1}_{-2}$。

后一个关系，是 $X_{0,1} = x^{-1}_{-2}$ 与 $X_{1,2} = x_{-2}^{-2} x_{-1} (1-x^{-1}_{-2})^{-2} = x^{-1}_{-3}$ 在 $\CK_{0,1}^2$ 之下的像。现在我们可以对一切 $n$ 继续用 $\CK_{n,n+1}$ 的逆作用，从而再算出一个无穷乘积。关系 $X_{n,n+1} = x^{-1}_{-n-2}$ 与 $X_{n-1,n} = x^{-1}_{-n-1}$ 在 $n\to\infty$ 时的极限为 $X_{1,1} = e^{-2a}$ 与 $X_{1,0} =- e^{-2 b}(1-e^{-2 a})^{-2}$。

于是，沿整条递推关系从 $n= - \infty$ 追到 $n=\infty$，我们便得到无穷乘积的一个表达式
\begin{equation} \label{eq:appA-infinite-product}
\cdots \CK_{4,3}^2 \CK_{3,2}^2 \CK_{2,1}^2 \CK_{1,0}^2 \CK_{0,1}^{-2}
\CK_{1,0}^{-2} \CK_{0,1}^2 \CK_{1,2}^2 \CK_{2,3}^2 \cdots
\end{equation}
递推关系两端极限值之间的映射为 $(X_{1,1},X_{1,0}) \to (X_{1,1}, X_{1,0}(1-e^{-2 a})^4 = X_{1,0}(1-e^{-4 a})^4 (1+e^{-2 a})^{-4})$，这正是所预期的 $\CK_{2,2}^2 \CK_{1,1}^{-4}$！

在正文中，我们把这一公式与具有两种味的 $SU(2)$ \SW 理论的穿墙行为联系起来（其味对称性为 $SO(4)=SU(2)_A \times SU(2)_B$）。我们曾论证，应当有一个类似的关系成立，它携带着关于味荷的信息。该关系应当给出乘积 $\CK_{1,0;1,0} \CK_{1,0;-1,0} \CK_{0,1;0,1} \CK_{0,1;0,-1}$ 的重排，也就是说，给出这样一个理论的穿墙公式：它有一个由电荷为 $(1,0)$ 的粒子组成的 $SU(2)_A$ 二重态（在 $SU(2)_A$ 与 $SU(2)_B$ 的嘉当生成元之下，味荷为 $(\pm 1,0)$），还有一个由电荷为 $(0,1)$ 的粒子组成的 $SU(2)_B$ 二重态（在 $SU(2)_A$ 与 $SU(2)_B$ 的嘉当生成元之下，味荷为 $(0,\pm 1)$）。

基本的变换为
\begin{equation} \label{eq:appA-flavor-A}
\CK_{1,0;1,0} \CK_{1,0;-1,0}: (X_{1,0},X_{0,1}) \to (X_{1,0},X_{0,1}(1-
\mu_A X_{1,0})^{-1}(1- \mu_A^{-1} X_{1,0})^{-1})
\end{equation}
以及
\begin{equation} \label{eq:appA-flavor-B}
\CK_{0,1;0,1} \CK_{0,1;0,-1}: (X_{1,0},X_{0,1}) \to (X_{1,0}(1- \mu_B
X_{0,1})(1- \mu_B^{-1} X_{0,1}),X_{0,1}).
\end{equation}

对这个问题，我们需要把来自 $SU(2)_A$ 或 $SU(2)_B$ 二重态中粒子的那些因子交替相乘。取
\begin{equation} \label{eq:appA-flavor-solution}
\begin{aligned}
& x_n = -\frac{1}{2}\frac{\cosh u \cosh v}{\sinh^2 a} + (-1)^n
\frac{1}{2}\frac{\sinh u \sinh v}{\sinh^2 a}  \\
& - \frac{\sqrt{(\cosh 2
a + \cosh 2 u)(\cosh 2 a + \cosh 2 v)}}{\sinh^2 2 a} \cosh (2 a n +
2 b)
\end{aligned}
\end{equation}
其中 $u,v$ 待下文用 $\mu_A, \mu_B$ 确定。

它满足递推关系
\begin{equation} \label{eq:appA-flavor-recursion}
x_{n+1} x_{n-1} = (1 - e^{u + (-1)^n v} x_n)(1 - e^{-u - (-1)^n v}
x_n).
\end{equation}

我们仍可令递推关系的初值为 $X_{0,1}=x_0$ 与 $X_{1,0}
= x^{-1}_1$。若施加 $\CK_{1,0;1,0} \CK_{1,0;-1,0}$，结果仍是 $X_{1,0}=x^{-1}_1$、$X_{2,1} =x^{-1}_2$，只要我们令 $\mu_A = e^{u-v}$。然后再施加 $\CK_{2,1;0,1}
\CK_{2,1;0,-1}$，结果便是 $X_{2,1}=x_2^{-1}$、$X_{3,2}= x^{-1}_3$，只要我们令 $\mu_B = e^{u+v}$。我们可以对一切 $n$ 交替地施加 $\CK_{2n+1,2n;1,0} \CK_{2n+1,2n;-1,0}$ 与 $\CK_{2n+2,2n+1;0,1} \CK_{2n+2,2n+1;0,-1}$，循着递推关系走到任意大的 $n$。借助关系 $X_{n+1,n} = x^{-1}_{n+1}$ 与 $X_{n,n-1} = x^{-1}_{n}$ 在 $n\to\infty$ 时的极限，便可算出这个无穷乘积。若取 $a,b$ 的实部为正，则 $X_{1,1} = e^{-2 a}$，且
\begin{equation} \label{eq:appA-flavor-forward}
X_{1,0} =- e^{-2 b} (1-e^{-4 a})^2(1+e^{-2a
-2u})^{-1/2}(1+e^{-2a +2u})^{-1/2}(1+e^{-2a -2v})^{-1/2}(1+e^{-2a
+2v})^{-1/2}.
\end{equation}

另一方面，我们也可以沿相反方向追随这个递推关系。关系 $X_{n,n+1} = x^{-1}_{-n-2}$ 与 $X_{n-1,n} = x^{-1}_{-n-1}$ 在 $n\to\infty$ 时的极限为 $X_{1,1} = e^{-2a}$，以及
\begin{equation} \label{eq:appA-flavor-backward}
X_{1,0} =- e^{-2 b} (1-e^{-4 a})^{-2}(1+e^{-2a
-2u})^{1/2}(1+e^{-2a +2u})^{1/2}(1+e^{-2a -2v})^{1/2}(1+e^{-2a
+2v})^{1/2}.
\end{equation}

总的映射为
\begin{equation} \label{eq:appA-flavor-total}
\begin{aligned}
& (X_{1,1},X_{1,0}) \to  \\
& \left(X_{1,1},X_{1,0} \frac{(1-X_{1,1}^2)^{4}}{(1+ \mu_A \mu_B X_{1,1})(1+ \mu_A^{-1} \mu_B X_{1,1})
(1+ \mu_A \mu_B^{-1}X_{1,1})(1+ \mu_A^{-1} \mu_B^{-1}X_{1,1})}\right).
\end{aligned}
\end{equation}

我们认出了所预期的答案：一个电荷为 $(2,2)$、不带味荷的矢量多重态，以及一个电荷为 $(1,1)$ 的超多重态，后者处于味对称性的 $(2_A)\otimes (2_B)$ 表示中（即 $SO(4)$ 的矢量）。

\begin{yizhu}
\begin{enumerate}[leftmargin=2em,itemsep=2pt,topsep=2pt]
\item \textbf{记号}：本附录沿用正文的记号。$\CK_{m,n}$ 是电荷为 $(m,n)$ 的 KS 辛同胚（即 KS 因子），$\CK_{m,n}^{2}$ 是这个因子的平方；带味荷时写作 $\CK_{m,n;s_1,s_2}$，分号后的 $(s_1,s_2)$ 是 $SU(2)_A\times SU(2)_B$ 的味荷。诸 $X_{m,n}$ 是相应电荷的达布坐标，例如本附录用到的 $X_{1,1}=X_{1,0}X_{0,1}$。
\item \textbf{勘误}：本附录中「然后再施加 $\CK_{2,1;0,1}\CK_{2,1;0,-1}$」一处，原文印作 $X_{2,1}={}^{-1}x_2$，上标位置有误，应为 $x_2^{-1}$；译文径改。
\end{enumerate}
\end{yizhu}
